% Appendix: Validation Burn-In Ablation

\subsection{Validation Burn-In Ablation (Figures~\ref{fig:val_burnin_ablation}--\ref{fig:val_burnin_budget})}
\label{appendix:val_burnin_ablation}

The paper's standard evaluation protocol runs the bandit through the
validation split ($n{=}1{,}785$) before evaluating on the held-out
test split ($n{=}1{,}824$).  To verify that this burn-in does not
artificially inflate test-set performance, we ablate the amount of
validation data used for online learning before test evaluation
\emph{across three budget regimes}: unconstrained, tight
($B{=}\$2.3{\times}10^{-4}$/req), and moderate
($B{=}\$6.6{\times}10^{-4}$/req).

\paragraph{Unconstrained regime.}
Five warmup conditions share identical priors (from train) and the
same hyperparameters ($\alpha{=}0.1$, $n_\mathit{eff}{=}1{,}000$,
$\gamma{=}0.995$); only the fraction of val used for burn-in varies
(0\%, 25\%, 50\%, 75\%, 100\%).  Two tabula-rasa baselines (0\% and
100\% burn-in, $n_\mathit{eff}{=}1$) complete a $2{\times}2$
factorial design crossing prior type with burn-in level.  All
conditions are evaluated on the same held-out test split across 20
random seeds with paired Wilcoxon signed-rank tests against the 100\%
burn-in reference.

\paragraph{Key findings (unconstrained).}
No warmup burn-in level differs significantly from the 100\% baseline
on test regret (Table~\ref{tab:val_burnin}).  The 0\% condition---warmup
priors with \emph{no} online learning before test---achieves $74.0$
cumulative regret, nominally \emph{lower} than the 100\% condition
($74.4$; $\Delta{=}{-}0.6\%$, $p{=}0.55$).  Early regret
(Regret@200) is lowest at 0\% burn-in ($5.9$ vs.\ $7.7$) because the
priors are fully intact at the start of the test trajectory.

\paragraph{Forgetting--burn-in interaction.}
With $\gamma{=}0.995$, effective memory is ${\sim}200$ steps.  After
1{,}785 val steps, warmup priors are decayed by $\gamma^{1785} \approx
1.3{\times}10^{-5}$---effectively erased.  This explains why the
warmup and tabula-rasa conditions converge after 100\% burn-in ($74.4$
vs.\ $73.9$; $p{=}0.67$): both operate on the same ${\sim}200$ recent
online observations, and the original priors contribute nothing.

\paragraph{Prior value: a deployment latency optimisation.}
The only statistically significant gap in the unconstrained regime is
between tabula rasa with \emph{no} burn-in ($83.5$) and all warmup
conditions ($74.0$--$76.0$; $p < 10^{-5}$), a $+12.2\%$ regret
penalty concentrated in the first ${\sim}200$ steps (R@200: $11.6$
vs.\ $5.9$).  Once sufficient online evidence accumulates---either
through val burn-in or directly on the test stream---tabula rasa
converges to the same steady-state (tabula rasa 100\% burn-in: $73.9$,
$p{=}0.67$ vs.\ warmup 100\%).

\paragraph{Budget-constrained regimes (Table~\ref{tab:val_burnin_budget}).}
To verify that these findings generalise to the budget-paced conditions
that are the paper's main contribution, we repeat the $2{\times}2$
factorial (priors~$\times$~burn-in at 0\%/100\%) under the Primal-Dual
BudgetPacer at tight and moderate targets.
To isolate the effect of reward-model learning from pacer calibration,
the BudgetPacer is \emph{pre-calibrated} on the full validation set:
all prompts are routed under the initial policy (priors or identity) and
the resulting costs are fed to the pacer, but the bandit is
\emph{not} updated.  All burn-in fractions therefore start the trial
with an identically calibrated $\lambda_t$; any remaining difference is
attributable solely to reward learning.%
\footnote{The oracle in this ablation is the quality-only oracle
$\max_a r_a$ (i.e., it ignores cost).  This inflates absolute regret
relative to the cost-adjusted oracle used in the main paper but does not
affect relative comparisons between conditions, which is the ablation's
purpose.}
The pattern is consistent across all regimes:
\begin{itemize}[nosep,leftmargin=*]
  \item \emph{Burn-in does not inflate.}  Within each budget regime,
    Warmup~0\% vs.\ Warmup~100\% is non-significant
    (tight: $198.2$ vs.\ $199.9$, $p{=}0.37$;
     moderate: $150.8$ vs.\ $149.4$, $p{=}0.70$).
  \item \emph{Priors matter at cold start (tight budgets).}
    Tabula Rasa with no burn-in has marginally higher regret at tight
    budgets ($205.1$ vs.\ $199.9$, $\Delta{=}{+}2.6\%$,
    $p{=}0.019$)---note this does not survive Bonferroni correction
    across the three within-regime tests.
    Prior to pacer pre-calibration this gap was larger
    ($\Delta{=}{+}6.0\%$, $p < 10^{-3}$), confirming that the original
    confound between cold-pacer and cold-bandit inflated the apparent
    harm of missing priors.
  \item \emph{Convergence after burn-in.}  With 100\% burn-in, warmup
    and tabula rasa converge at every budget level
    (tight: $199.9$ vs.\ $197.2$, $p{=}0.11$;
     moderate: $149.4$ vs.\ $144.4$, $p{=}0.28$),
    confirming that geometric forgetting erases the prior regardless of
    budget pressure.
  \item \emph{Moderate budgets collapse the gap.}  At moderate budgets,
    even tabula rasa without burn-in is non-significant vs.\ warmup
    ($152.2$ vs.\ $149.4$, $p{=}0.29$): the budget constraint
    collapses the feasible arm set, reducing the harm of uninformed
    exploration.
  \item \emph{Budget compliance is consistent.}  Cost/target ratios
    range from $97.6\%$ to $100.0\%$ across all conditions, confirming
    that the BudgetPacer achieves comparable spending regardless of
    burn-in level or prior initialisation.
\end{itemize}

\paragraph{Implication.}
Warmup priors are best understood as a \emph{deployment latency
optimisation}: they eliminate the cold-start transient where uninformed
exploration dominates, halving early regret (R@200: $5.9$ vs.\ $11.6$
unconstrained; $23.1$ vs.\ $27.6$ at tight budget).
Their contribution is transient by design---geometric forgetting
replaces offline estimates with online evidence within ${\sim}200$
effective-memory steps---so the router's steady-state quality is
determined by the online learner and budget pacer, not the
initialisation.
This holds \emph{regardless of budget regime}:
the BudgetPacer's adaptive $\lambda_t$ provides the same closed-loop
cost control whether the router starts from informed priors or from
scratch.
The case for warmup priors thus scales with the deployment's sensitivity
to early-stage regret: applications routing thousands of requests in
the first minutes benefit substantially; systems that can tolerate a
brief learning phase may forgo priors entirely and rely on the online
learner alone.

\begin{table}[t]
\centering
\caption{Validation burn-in ablation---unconstrained regime (20~seeds,
held-out test split).
$\Delta$ is relative to the 100\% burn-in reference.  No warmup
condition differs significantly from the reference; the only
significant gap is tabula rasa without burn-in.}
\label{tab:val_burnin}
\small
\begin{tabular}{@{}llrrrl@{}}
\toprule
\textbf{Priors} & \textbf{Burn-in} & \textbf{Regret} & \textbf{R@200} & $\boldsymbol{\Delta}$ & $\boldsymbol{p}$ \\
\midrule
Warmup & 0\%   & 74.0\,$\pm$\,0.5 & 5.9  & $-0.6$\% & 0.55 \\
Warmup & 25\%  & 76.0\,$\pm$\,0.7 & 8.6  & $+2.2$\% & 0.22 \\
Warmup & 50\%  & 75.1\,$\pm$\,0.8 & 7.9  & $+0.9$\% & 0.50 \\
Warmup & 75\%  & 75.1\,$\pm$\,0.6 & 7.9  & $+1.0$\% & 0.52 \\
Warmup & 100\% & 74.4\,$\pm$\,0.8 & 7.7  & ---       & --- \\
\midrule
Tabula Rasa & 0\%   & 83.5\,$\pm$\,0.8 & 11.6 & $+12.2$\% & $< 10^{-5}$ \\
Tabula Rasa & 100\% & 73.9\,$\pm$\,0.7 & 8.2  & $-0.8$\%  & 0.67 \\
\bottomrule
\end{tabular}
\end{table}

\begin{table}[t]
\centering
\caption{Validation burn-in ablation---budget-constrained regimes
(20~seeds, held-out test split, BudgetPacer active, pacer pre-calibrated).
$\Delta$ is relative to the Warmup~100\% reference within each regime.
\emph{Cost/T} is the mean cost per request as a fraction of the budget
target (1.00~=~at target).
The burn-in finding holds under budget pressure: within each regime,
Warmup~0\% is non-significant vs.\ Warmup~100\%, and cost compliance is
consistent across conditions.}
\label{tab:val_burnin_budget}
\small
\begin{tabular}{@{}lllrrrrl@{}}
\toprule
\textbf{Budget} & \textbf{Priors} & \textbf{Burn-in} & \textbf{Regret} & \textbf{R@200} & \textbf{Cost/T} & $\boldsymbol{\Delta}$ & $\boldsymbol{p}$ \\
\midrule
Tight    & Warmup      & 0\%   & 198.2\,$\pm$\,1.0 & 23.1 & 1.00 & $-0.8$\%  & 0.37 \\
         & Warmup      & 100\% & 199.9\,$\pm$\,1.3 & 23.0 & 1.00 & ---        & --- \\
         & Tabula Rasa & 0\%   & 205.1\,$\pm$\,2.1 & 27.6 & 0.99 & $+2.6$\%  & 0.019 \\
         & Tabula Rasa & 100\% & 197.2\,$\pm$\,1.1 & 21.6 & 1.00 & $-1.3$\%  & 0.11 \\
\midrule
Moderate & Warmup      & 0\%   & 150.8\,$\pm$\,3.0 & 18.3 & 0.99 & $+0.9$\%  & 0.70 \\
         & Warmup      & 100\% & 149.4\,$\pm$\,1.9 & 16.5 & 0.99 & ---        & --- \\
         & Tabula Rasa & 0\%   & 152.2\,$\pm$\,3.0 & 20.8 & 0.98 & $+1.9$\%  & 0.29 \\
         & Tabula Rasa & 100\% & 144.4\,$\pm$\,2.1 & 16.0 & 0.99 & $-3.4$\%  & 0.28 \\
\bottomrule
\end{tabular}
\end{table}
